\chapter*{Evaluating Quality of Receptors for Molecular Docking}

{
\color{red}{It is possible on many occasions that several entries for the same receptor are available on the wwPDB. One must choose based on the conformational state of the receptor that is desired. Then selecting the entry with a good resolution, in our opinion, a resolution below 2.0 \r{A} is a good choice whenever possible. There are other structure quality parameters, apart from the resolution, that provide useful information about the target. the b-factors or beta-factors provide key information on the flexibility of the atoms. It also holds key information on the confidence with the crystallographers are able to narrow down the atomic positions while fitting the electron density maps, for examples atoms with b-factors $\le40 $\r{A}$^2$$ $ provide good confidence in their assignments. While using the CSMs, predicted local distance difference test (pLDDT), ranges from 0 to 100, depicts a per-residue measure of local confidence that must be carefully examined. pLDDT values $\ge70$ suggest good confidence in the prediction, while values $<50$ indicate very low confidence. The users are directed to the \href{https://pdb101.rcsb.org/learn/guide-to-understanding-pdb-data/computed-structure-models}{\textit{Guide to Understanding PDB Data}} with a section dedicated to CSMs  The choice of receptor should also bear in mind if there exists an entry with a co-crystallised ligand in the binding site of interest.  While preparing the receptor, the atomic resolution limitation makes it difficult to locate the position of hydrogens, which must be added in the preparation step. Partial atomic charges are calculated and assigned for individual atoms.  The decision to keep important metals, co-factors, and water molecules in the docking site is made at this stage. If they are crucial for ligand binding, they must be retained; otherwise, they can be removed. Such decisions are far more complex than they appear at first sight. To come to a meaningful decision, the computational chemist must be well aware of the biology and functioning of the receptor in question (\textbf{see  Tip 1}). Recently, there has been a particular interest in developing tools that systematic evaluate the energetics of displacing water molecules from the binding site by the ligands. In doing so, they are well placed to predict crucial water molecules that may lead to significant entropic gains. The tools such as  waterMap\cite{bib48, bib49}, WaterDock\cite{bib47} and SSTmap\cite{bib50} are available that predicts the energetics of displacing water molecules from the binding site. Deleting co-factors, waters, or metal ions deemed important can lead to meaningless results, even though very high or favourable docking scores are observed.

\subsubsection*{File formats for macromolecules:}
    \begin{enumerate}
        \item \textit{PDB (Protein Data Bank):} Stores 3-Dimensional structural data of proteins and nucleic acids. 
	       \item \textit{PQR:} A variant of PDB with atomic charges and radii, often used in electrostatics calculations.
	       \item \textit{MCIF/MMTF:} Compressed formats for macromolecular structures, useful for large datasets. 
    \end{enumerate}


Generally, it is highly recommended to perform energy minimization using a force field compatible with the docking program. This step optimizes all the inconsistencies in the structure or shows warnings on the errors that can be detrimental to the reliability and accuracy of the results. 
}}

